\BourbakiExerciseGroup

\begin{enumerate}
\def\labelenumi{\arabic{enumi})}
\tightlist
\item
  一个拓扑 \(\mathfrak{G}\) （在有序交换群 \(\mathbf{G}\) 上）称为与 \(\mathbf{G}\) 的序群结构相容，如果它与 \(\mathbf{G}\) 的群结构相容，此外集合 \(\mathbf{G}_{+}\) （由元素 \(x \geqslant 0\) （属于 \(\mathbf{G}\) ）组成）关于 \(\mathfrak{G}\) 是闭的。
\end{enumerate}

\begin{enumerate}
\def\labelenumi{\alph{enumi})}
\item
  设 \(\mathfrak{G}\) 是与 \(\mathbf{G}\) 的序群结构相容的一个 Hausdorff 拓扑，又设 \(\hat{\mathbf{G}}\) 是 \(\mathbf{G}\) 的完备化（关于 \(\mathfrak{G}\) ）。在群 \(\hat{\mathbf{G}}\) 中，闭包 \(\mathbf{P}\) （\(\mathbf{G}_{+}\) 的）是这样的：\(y - x \in \mathbf{P}\) 是一个与 \(\hat{\mathbf{G}}\) 的群结构相容的预序。
\item
  设 \(\theta\) 是一个无理数，考虑 \(\mathbf{Q}^2\) 上的这样一个序：其正元素组成的集合由 (0, 0) 与所有偶 \((x, y)\) （满足 \(y - \theta x \geqslant 0\) ）组成。关于这个序，\(\mathbf{Q}^2\) 是一个全序群，而 \(\mathbf{Q}^2\) 上由有理直线的拓扑与自身的乘积给出的拓扑与这个序群结构相容。但在 \(\hat{\mathbf{G}} = \mathbb{R}^2\) 上，关系 \(y - x \in \mathbf{P} = \overline{\mathbf{G}}_+\) 不是序。
\item
  在全序群 G 上，拓扑 \(\mathfrak{C}_{0}(\mathbf{G})\) （第一章，§ 2，习题 5）是与序群结构相容的所有拓扑中最粗的（区分两种情形，视所有 \(x > 0\) 组成的集合有还是没有最小元素而定）。G 上与 G 的群结构相容且细于 \(\mathfrak{C}_{0}(\mathbf{G})\) 的每个拓扑都与 G 的序群结构相容，且映射 \(x \to x^{+}\) 关于这样的拓扑连续。这个映射关于 \(\mathfrak{C}_{0}(\mathbf{G})\) 一致连续，但不一定关于一个群拓扑 \(\mathfrak{C}\) （严格细于 \(\mathfrak{C}_{0}(\mathbf{G})\) ）一致连续{[}见 b){]}。
\end{enumerate}

\begin{enumerate}
\def\labelenumi{\arabic{enumi})}
\setcounter{enumi}{1}
\tightlist
\item
  \begin{enumerate}
  \def\labelenumii{\alph{enumii})}
  \tightlist
  \item
    设 G 是一个格序交换群。对 G 中每个 \(x \geqslant 0\) ，令 \(\mathbf{I}(x)\) 表示 G 中的区间 \([-x, x]\) 。G 的一个非空族 \((c_{\alpha})\) （由元素 \(\geqslant 0\) 组成）使得诸区间 \(\mathbf{I}(c_{\alpha})\)
  \end{enumerate}
\end{enumerate}

构成 G 上与 G 的群结构相容的拓扑 \(\mathfrak{G}\) 在点 0 处的基本邻域系，当且仅当 \(c_{\alpha}\) 组成的集合是有向的（关于关系 \(\geqslant\) ），且对每个 \(\alpha\) 存在 \(\beta\) 使 \(2c_{\beta} \leqslant c_{\alpha}\) 。如果这些条件满足，G 到 G 中的映射 \(x \to x^{+}\) 一致连续。\(\mathfrak{G}\) 是 Hausdorff 的当且仅当 \(\inf c_{\alpha} = 0\) ；且在这种情形 \(\mathfrak{G}\) 与 G 的序群结构相容。

\begin{enumerate}
\def\labelenumi{\alph{enumi})}
\setcounter{enumi}{1}
\item
  在群 \(G = Z^{N}\) （可数无穷多个全序群 Z 的乘积）上，没有用 a) 的程序定义的非离散拓扑，但诸因子 Z 上的离散拓扑的乘积是 G 上的一个拓扑，它与 G 的序群结构相容，它是 Hausdorff 的且非离散的。
\item
  在全序群 G 上，用 a) 的程序定义的唯一 Hausdorff 拓扑是拓扑 \(\mathcal{G}_{0}(G)\) （第一章，§ 2，习题 5）。
\end{enumerate}

\begin{enumerate}
\def\labelenumi{\arabic{enumi})}
\setcounter{enumi}{2}
\tightlist
\item
  设 G 是一个格序交换群，又设 M 是 G 的上界子集组成的半群。M 的每个上有界（相应地，下有界）子集都有上确界（相应地，下确界），且 G 可以按典型方式与 M 的一个子集等同。设 G' 表示 M 的最大子群（M 的所有可对称化元素组成的集合）。
\end{enumerate}

\begin{enumerate}
\def\labelenumi{\alph{enumi})}
\tightlist
\item
  考虑一个拓扑 \(\mathcal{G}\) （在 \(\mathbf{G}\) 上，由习题 2 \(a\) ) 的程序定义）；设 \(\mathbf{V}_{\alpha}\) 是所有偶 \((z, z')\) （属于 \(\mathbf{M} \times \mathbf{M}\) ）组成的集合，使得
\end{enumerate}

\[
z - c _ {\alpha} \leqslant z ^ {\prime} \leqslant z + c _ {\alpha};
\]

证明诸 \(V_{\alpha}\) 构成 M 上一个一致结构 U 的一致邻域基本系，且 U 诱导的拓扑 \(C'\) 在 G 上诱导出 C。若 C 是 Hausdorff 的，则 U 是一个 Hausdorff 一致结构（注意 M 的元素——看作 G 的子集——这时是闭的），且这样定义的一致空间 M 是完备的。证明 M × M 到 M 中的映射 \((z, z') \to z + z'\) 一致连续。假设 C 是 Hausdorff 的，证明在 M 中 M 的任意子集的上（相应地，下）界组成的集合关于 \(C'\) 是闭的（先对 M 的一个元素的上界或下界组成的集合证明此点，把 M 的元素看作 D 的子集），且映射

\[
(z, z ^ {\prime}) \rightarrow \inf (z, z ^ {\prime})
\]

\(\mathbf{M} \times \mathbf{M}\) 到 \(\mathbf{M}\) 中的这一映射是一致连续的（同样的方法）。

\begin{enumerate}
\def\labelenumi{\alph{enumi})}
\setcounter{enumi}{1}
\item
  在本习题中从现在起设 \(\mathfrak{G}\) 是 Hausdorff 的。这时闭包 \(\overline{\mathbf{G}}\) （\(\mathbf{G}\) 在 \(\mathbf{M}\) 中的）是一个格序群，它在 \(\mathfrak{G}'\) 诱导的拓扑中是完备的。群 \(\mathbf{G}'\) 在 \(\mathbf{M}\) 中是闭的（注意它的闭包是 M 的一个子群），且 \(\mathcal{G}'\) 在 \(\mathbf{G}'\) 上诱导的拓扑与 \(\mathbf{G}'\) 的序群结构相容。为使 \(\mathbf{G}'\) 等于 \(\overline{\mathbf{G}}\) ，只需每个非空开区间 \(]a, b[\) （\(\mathbf{G}'\) 中的）含有 \(\mathbf{G}\) 的一个元素，且若 \(\mathbf{G}\) 是全序的这总是成立的。在这种情形，\(\mathbf{G}'\) 上由 \(\mathcal{G}'\) 诱导的拓扑是 \(\mathcal{C}_0(\mathbf{G}')\) 。
\item
  取 G 为有序群 \(Q^{2}\) ，即全序群 Q 与自身的乘积；这时带乘积序有 \(G' = M = R^{2}\) 。证明 G 上由习题 2 a) 的程序定义的非离散 Hausdorff 拓扑只有三个，且它们是这样得到的：取元素 \(c_{\alpha}\) 组成的集合为下列三者之一：偶 \((x, y)\) （满足 x \textgreater{} 0 与 y \textgreater{} 0 ）组成的集合，偶 \((x, 0)\) （满足 x \textgreater{} 0 ）组成的集合，或偶 \((0, y)\) （满足 y \textgreater{} 0 ）组成的集合。对这三个拓扑中的第一个有 \(\overline{G} = G'\) ，尽管 \(G'\) 有不含 G 的任何元素的非空开区间；对另两个拓扑有 \(\overline{G} \neq G'\) 。对这三个拓扑的每一个，G 中元素 z \textgreater{} 0 组成的集合不是开的。
\item
  取 G 为赋以字典序的群 \(Q^{2}\) （集合论，第三章，§ 2，第 6 节）；G 是一个非阿基米德全序群。这时 \(\overline{G}=G'\neq M\) ；M 是全序的，但 M 上的拓扑 \(C'\) 不同于拓扑 \(\mathcal{C}_{0}(M)\) ；\(G'\) 同构于赋以字典序的群 \(R\times Q\) ；在 \(G'\) 中子群 \(H=R\times\{0\}\) 是开的且孤立的，但在商群 \(G'/H\) 上商拓扑不同于 \(\mathcal{C}_{0}(G'/H)\) 。
\end{enumerate}

4) a) 设 E 是一个全序集。设 F 是 E 的元素以 Z 为系数的形式有限线性组合组成的 Z 模，并在 F 上定义一个全序群结构，取集合 \(F_{+}\) （由元素 \(\geqslant 0\) 组成）为由 0 与所有线性组合 \(\sum_{\xi \in E} n(\xi) \xi \neq 0\) （它们满足 \(n(\xi) > 0\) ，对最大元素 \(\xi \in E\) （使 \(n(\xi) \neq 0\) ）而言）组成的集合。证明 E 可以与 \(F_{+}\) 的一个共尾子集等同。

\begin{enumerate}
\def\labelenumi{\alph{enumi})}
\setcounter{enumi}{1}
\item
  设 E 是良序的，考虑 E 到 F 中的有界映射组成的群 G。在 G 上定义一个全序群结构，取集合 \(G_{+}\) （由元素 \(\geqslant0\) 组成）为由 0 与所有 E 到 F 中的有界映射 x（它们满足 \(x(\xi)>0\) ，对最小 \(\xi\in E\) （使 \(x(\xi)\neq0\) ）而言）组成的集合（字典序）。证明 G 关于拓扑 \(\mathcal{C}_{0}(G)\) 是完备的；此外若 E 不可数且每个截段 \(\leftarrow,\xi]\) ( \(\xi\in E\) ) 是可数的，则 G 中开集的每个可数交是开的，且 G 的每个紧子集是有限的（cf.~第一章，§ 9，习题 4）。
\item
  此后设 E 是良序的且不可数，但每个截段 \(\leftarrow,\xi]\) 是可数的。对每个 \(\xi\in E\) ，令 \(c_{\xi}\) 是 E 到 E 中的常值映射，它在每点的值为 \(\xi\) 。诸 \(c_{\xi}\) 在 \(G_{+}\) 中构成一个共尾子集。设 \(E_{0}\) 是由 E 添加一个元素 \(\omega\) 而得到的集合，考虑集合 \(X=G\times E_{0}\) 上由下列集合生成的拓扑 G：(i) \(V_{a,b,\xi} = ]a\) , \(b[\times\{\xi\}\) ，其中 G 中 a \textless{} b 且 \(\xi\in E\) ；(ii) \(W_{a,b,\xi} -\) ，对 G 中 a \textless{} b 且 \(\xi\in E\) 使 \(c_{\xi} > \sup(|a|, |b|) -\) ——由偶 \((x, \omega)\) （满足 a \textless{} x \textless{} b ），以及偶 \((x, \zeta)\) （满足 \(\zeta\in E\) 、\(\zeta \geqslant \xi\) 且或 a \textless{} x \textless{} b 或 \(4c_{\zeta} + a < x < 4c_{\zeta} + b\) ）组成。证明 G 是 Hausdorff 的，且 X 的每个紧子集（关于 G）是有限的。此外 G 依法则 \((x, (y, \zeta)) \to (x + y, \zeta)\) 连续地作用在 X 上；但 G 不逆紧地作用于 X，尽管第三章 § 4 第 2 节命题 4 的条件 a)、b)、c)、d) 满足，且对 X 的每对紧子集 K, L，集合 P(K, L)（第三章，§ 4，第 5 节，定理 1）是紧的。
\end{enumerate}
% semantic-fragments: legacy-input-seam
\relax{}
